\documentclass{article}
\usepackage[T1]{fontenc}
\usepackage{lmodern}
\usepackage{graphicx}
\usepackage{amsmath,amssymb}
\usepackage[a4paper,margin=2cm]{geometry}
\usepackage[absolute,overlay]{textpos}
\setlength{\TPHorizModule}{1mm}
\setlength{\TPVertModule}{1mm}
\usepackage[indonesian]{babel}
\usepackage{hyperref}
\setlength{\emergencystretch}{2em}
% unit: Halaman 9

\begin{document}

% === Halaman 9 (python/native) ===

• Kurva eliptik y2 = x3 + ax + b terdefinisi untuk x,y  R

• Didefinisikan sebuah titik bernama titik O(x, ), yaitu titik pada 

infinity.

• Titik-titik P(x, y) pada kurva eliptik bersama operasi + membentuk 

sebuah grup.


Himpunan G      : semua titik P(x, y) pada kurva eliptik


Operasi biner   : +

• Penjelasan kenapa kurva eliptik membentuk sebuah grup dijelaskan 

pada slide-slide berikut ini.

9

\end{document}
